\section{Discussion}
\label{sec:discussion}

\subsection{Why simple rules suffice}
\label{sec:disc-limits-protocol}
Three results bound what a more elaborate reporting protocol could add under the gaps measured here. First, the timing of recovery probes is near its limit: a schedule fitted to each record's gap lengths saved at most \probeSaveLoraMax{} of the design's energy on LoRaWAN (Section~\ref{sec:hand-batch}); only on NB-IoT, where every failed probe pays an unmeasured cell search, did it reach \probeSaveNbMax. Second, the reporting decisions act on a shrinking share of the budget: with hourly batching, transmission was \txShareFinalMin--\txShareFinalMax{} of the node's energy; therefore, even a protocol that halved it would extend battery life far less than the transmission ratios suggest. Third, the parts of the design are separable by what they buy: the bounded-error trigger buys the real-time picture, the alarm upload buys alarm latency, and the hourly batch buys completeness within the hour at \noBatchSaveMin--\noBatchSaveMax{} of the transmission energy; an operator who does not need the last can drop it. The remaining gains lie outside the reporting protocol: in sensing time and sensor power, in the network that causes the gaps, and in measuring the target hardware.

\subsection{Practical guidance}
\label{sec:disc-guidance}
The results lead to six recommendations:
\begin{enumerate}
\item Determine which regime loses the data: Frame counters, a log of gateway and server downtime, and the node's own status are sufficient to distinguish short radio losses from long delivery gaps, which call for different remedies; models that assume geometric gap lengths should be checked against such a log.
\item Do not fight long gaps with retransmissions: Confirm, buffer, and back off; on NB-IoT, back off even without a buffer, because searching for a cell in every slot of a gap cost up to \backoffMax{} the energy of backing off. Give the buffer tens of kilobytes, and at SF12 let live readings go ahead of the backlog; where many nodes share a gateway, randomise recovery probes.
\item Batch, and let the batch leave early when it matters: Hourly batching delivers every reading for a fraction of per-reading energy at SF7 and on NB-IoT; an alarm trigger restores timeliness, and a bounded-error trigger that mirrors the server's estimate restores the real-time picture, with the tolerance as the operator's energy control. The node must model when its data arrive and, on LoRaWAN, whether the application has them. Where the record need not be complete within the hour, the trigger without the batch saves a further \noBatchSaveMin--\noBatchSaveMax{} of transmission energy.
\item Count the whole node: Once readings are batched, transmission is only \txShareFinalMin--\txShareFinalMax{} of the node's energy; therefore, the probe warm-up time and the sensor board's power deserve as much attention as the radio protocol.
\item Judge reconstruction against the site's own diel pattern, the best simple real-time estimator for DO and pH: A test against the last value alone overstates what neighbours and models add.
\item Engineer the gaps themselves: No reporting scheme shortens a gap. Gateway diversity masked a month-long gateway outage at UVA, and watchdog resets bound silent gateway failures~\cite{brzozowski2018lessons,ingelrest2010sensorscope}.
\end{enumerate}

\subsection{Validity and limitations}
\label{sec:disc-limits}
The main limitations are as follows:
\begin{itemize}
\item Records: The loss evidence comes from four records and supports statements about them, not a law for all deployments; Weija serves as qualitative support, and the cause of a gap could be checked only partly; therefore, buffering's recovery is an upper bound.
\item Energy without hardware: No hardware was built for this study. The model rests on published measurements of specific modules, and an independent measurement on another board differs severalfold in absolute terms. The claims are therefore ratios on the same radio, tested against a factor-of-two uncertainty in every component (Section~\ref{sec:hand-network}). At SF12, the advantage of batching, and everywhere, the sign of buffering's small energy difference, need hardware measurement, which is the most important next step.
\item Model scope: Clock faults and storage failure are not simulated, and the multi-node model of one gateway shows mechanisms (synchronised probes, acknowledgement capacity) rather than numbers for a particular network.
\item Events and rivers: Events are synthetic steps and ramps of realistic size, and the detector catches slow ramps poorly (\detectRampFive{} of 5$\sigma$ ramps); the latency comparisons concern detected events. The river data cover two years in one country, and the held-out year shares most sites with the development year (\sitesBothYears{} of \confSites); therefore, it tests the design on new data rather than on new rivers.
\item Reconstruction: The deep imputers were trained on the first half of one year per river chain, which is a small amount of data for such models, and the neighbours lie \pairKmMin{} to \pairKmMax~km from their targets (median \pairKmMed~km); denser networks and longer training records may reconstruct gaps better.
\item Measurement resolution and pre-specification: pH is recorded in steps of 0.1, equal to its tolerance at the reference level; therefore, part of the trigger's activity and of the picture score for pH responds to one-step flicker rather than to changes in the river. Hypotheses were written before their runs, and the confirmatory code was frozen before the held-out data were downloaded, but neither was registered with a third party; every verdict is reported, including the failures and the defect (D5) found before the held-out run.
\end{itemize}

\subsection{Conclusion}
In the four records studied, data were lost in two regimes: short radio losses and rarer long delivery gaps that, in the water-quality deployments, held most of the missing data. Retransmission spends energy on gaps it cannot bridge, and the evaluated reconstruction estimators, simple or deep, kept most day-long gaps outside the tolerance. Buffering with back-off delivers every generated reading at an energy \bufLoMin{} to \bufLoMax\% different from that of not recovering at all on LoRaWAN, and hourly batching makes it cheaper still. A design that batches hourly, sends alarms immediately, and leaves early whenever the server's estimate would leave tolerance matched per-reading reporting on alarm latency and on the real-time picture on a held-out year, while per-reading ARQ used \arqOverBoundedMinConf--\arqOverBoundedMaxConf{} as much transmission energy and \nodeArqOverFinalMin--\nodeArqOverFinalMax{} as much node energy. A better probing schedule would save little, and once readings are batched, transmission is at most approximately half of the node's energy. What remains is the delivery gap itself, which can only be shortened by network engineering. Future work includes hardware measurements, longer delivery records from tropical deployments, the energy of NB-IoT cell search without a network, and recovery across several gateways.
